\begin{center}
    {\Large\bfseries 4. ALCANCE DE LA INVESTIGACIÓN}
\end{center}

Este estudio es de tipo explicativo: busca establecer si la integración e interpretación contextualizada de métricas avanzadas de rendimiento accesible (variable independiente / causa) produce una diferencia medible en la precisión de la caracterización funcional y en la reducción de la incertidumbre al diferenciar perfiles de atacantes (variables dependientes / efectos), en comparación con un modelo de evaluación tradicional basado en estadísticas brutas (goles y asistencias). La elección de este alcance se justifica en que la investigación no se limita a describir el uso de datos en el fútbol ni a observar la simple correlación entre variables, sino que pone a prueba el impacto de un marco de ponderación táctico-contextual sobre la evaluación de jugadores. Para ello, se apoya en el diagnóstico descriptivo previo sobre la toma de decisiones en el *scouting* (numeral 1) y en los estudios correlacionales existentes sobre la efectividad de métricas como goles esperados ($xG$) y asistencias esperadas ($xA$) en clubes de élite, evaluando si esta metodología permite diferenciar objetivamente a futbolistas con registros tradicionales similares en la Premier League durante la temporada 2025/2026.

Limitaciones de la investigación

Al basarse en datos secundarios de eventos (*event data*) de acceso abierto y no en un entorno controlado con datos de seguimiento óptico en tiempo real (*optical tracking data*), las inferencias sobre comportamientos tácticos sin balón deben asumirse con mayor cautela que las derivadas de mediciones de alta resolución de clubes de élite; la muestra se limita de manera exclusiva a futbolistas de perfil ofensivo con minutos jugados en la Premier League durante la temporada 2025/2026, lo que restringe qué tanto se pueden generalizar o transferir directamente las ponderaciones cuantitativas a otras ligas continentales, divisiones inferiores o demarcaciones defensivas sin una recalibración previa de los parámetros competitivos; y el modelo de ponderación multicriterio diseñado para matizar las métricas según el contexto del equipo se valida dentro de este mismo estudio, por lo que sus propiedades metodológicas (consistencia interna, sensibilidad del modelo y nivel de concordancia) deben reportarse empíricamente junto con los resultados finales del proyecto, no darse por supuestas previamente.

